\documentclass[11pt]{amsart}

% Encoding and fonts
\usepackage[T1]{fontenc}
\usepackage[utf8]{inputenc}
\usepackage{lmodern}

% Math packages
\usepackage{amsmath}
\usepackage{amssymb}
\usepackage{amsfonts}
\usepackage{amsthm}
\usepackage{mathtools}

% Layout and utilities
\usepackage{geometry}
\usepackage{enumitem}
\usepackage[hidelinks]{hyperref}

\geometry{margin=1in}

\numberwithin{equation}{section}
\setcounter{tocdepth}{3}
\setcounter{secnumdepth}{3}

% ---------------------------------------------------------
% Theorem-like environments
% ---------------------------------------------------------

\theoremstyle{plain}
\newtheorem{theorem}{Theorem}[section]
\newtheorem{proposition}[theorem]{Proposition}
\newtheorem{lemma}[theorem]{Lemma}
\newtheorem{corollary}[theorem]{Corollary}
\newtheorem{conjecture}[theorem]{Conjecture}

\theoremstyle{definition}
\newtheorem{definition}[theorem]{Definition}
\newtheorem{experiment}[theorem]{Experiment}
\newtheorem{heuristic}[theorem]{Heuristic}
\newtheorem{question}[theorem]{Question}

\theoremstyle{remark}
\newtheorem{remark}[theorem]{Remark}
\newtheorem{warning}[theorem]{Warning}

\newcommand{\N}{\mathbb{N}}
\newcommand{\Z}{\mathbb{Z}}
\newcommand{\C}{\mathbb{C}}
\newcommand{\Hil}{\mathcal{H}}

\title{
Notebook:
Non-Commutative and Non-Hilbert Structures
Attached to the Ordinal--M\"obius Pairing
}

\author{}
\date{}

\begin{document}

\maketitle

\tableofcontents

\bigskip

\begin{remark}
\label{rem:exploratory_nature}
This notebook is exploratory in nature. Its purpose is to investigate
whether the pair $(n,\mu(n))$ -- the ordinal index and the value of
the M\"obius function at that index -- can be organized into a
coherent operator-theoretic structure exhibiting a genuine,
non-statistical incompatibility analogous to the position-momentum
complementarity of quantum mechanics.

Several sections contain formal manipulations, heuristic arguments,
and speculative interpretations. Whenever a statement is rigorous it
is explicitly identified as such (Definition, Proposition, Lemma,
Corollary), with a full derivation. All remaining claims are
tentative, heuristic, or conjectural, and are marked as such.

The scope of this notebook is the algebraic and operator-theoretic
structure of the ordinal--M\"obius pairing itself. It does not
address the asymptotic behaviour of the Mertens function $M(x)=\sum_{k\le x}\mu(k)$
or the Riemann Hypothesis; no result below bears on either.
\end{remark}

\section{Introduction}
\label{sec:introduction}

The M\"obius function $\mu:\N\to\{-1,0,1\}$ is defined by $\mu(1)=1$,
$\mu(n)=(-1)^{\omega(n)}$ if $n$ is squarefree with $\omega(n)$ distinct
prime factors, and $\mu(n)=0$ otherwise. The present notebook records
an attempt to attach to the pair $(n,\mu(n))$ an operator-theoretic
structure in which $n$ (the ordinal) plays the role of a
\emph{position} variable and $\mu(n)$ (or a quantity built from it)
plays the role of a \emph{momentum}-type variable, in such a way that
knowledge of one genuinely precludes exact knowledge of the other --
not as a statement about measurement precision, but as a structural
fact about the objects involved.

The motivating question, stated informally, is:

\begin{question}
\label{q:motivating}
Can the ordinal $n$ and the M\"obius value $\mu(n)$ be realized as a
canonically conjugate pair, in a sense that reproduces the qualitative
content of Heisenberg complementarity (an intrinsic impossibility, not
a statistical spread), while remaining fully precise and auditable at
each step?
\end{question}

Section~\ref{sec:naive_ansatz} records a first, heuristic attempt
(``MQMNL'') and explains why it does not answer
Question~\ref{q:motivating}: it postulates an uncertainty relation
without deriving it from the algebra of the operators involved.
Section~\ref{sec:obstruction} isolates a precise obstruction to the
most literal reading of the question. Section~\ref{sec:number_phase}
gives a Hilbert-space construction, adapted from the number-phase
formalism of quantum optics, in which a genuine (derived, not
postulated) complementarity holds between the ordinal and a phase
variable, with $\mu$ recovered as a coarse-grained observable.
Section~\ref{sec:non_hilbert} gives a second, non-Hilbert-space
construction, built directly from the Dirichlet convolution ring of
arithmetic functions, in which the incompatibility between $n$ and
$\mu$ is an algebraic fact with no probabilistic content at all.
Section~\ref{sec:literature} discusses the closest structures found in
the literature. Section~\ref{sec:open} lists open questions.

\section{A naive quantum ansatz, and why it fails}
\label{sec:naive_ansatz}

\begin{remark}[Failed attempt]
\label{rem:failed_mqmnl}
An initial heuristic proposal (``M\"obius Quantum Model, Non-Localizable'',
MQMNL) represents the partial sum $M(n)=\sum_{k\le n}\mu(k)$ as a state
$|\Psi_n\rangle=\sum_{m\in\Z}a_m^{(n)}|m\rangle$ in $\ell^2(\Z)$,
evolving by
\begin{equation}
|\Psi_n\rangle = T_{\mu(n)}|\Psi_{n-1}\rangle,
\qquad
T_{+1}|m\rangle=|m+1\rangle,\
T_{0}=\mathrm{Id},\
T_{-1}|m\rangle=|m-1\rangle,
\label{eq:mqmnl_evolution}
\end{equation}
and postulates an uncertainty relation
\begin{equation}
\Delta\mathcal{M}_n \, D_{\theta,n} \geq \kappa
\label{eq:mqmnl_postulated_uncertainty}
\end{equation}
for some constant $\kappa>0$, where $D_\theta$ is a dispersion of a
conjugate phase variable.

This does not answer Question~\ref{q:motivating}, for an identifiable
reason: the operators $T_{+1},T_0,T_{-1}$ in \eqref{eq:mqmnl_evolution}
all commute with one another, being translations on $\Z$, which forms
an abelian group. A translation does not broaden a distribution, so
the evolution \eqref{eq:mqmnl_evolution} alone provides no mechanism
that could force \eqref{eq:mqmnl_postulated_uncertainty}. The relation
\eqref{eq:mqmnl_postulated_uncertainty} is an additional hypothesis,
not a consequence of \eqref{eq:mqmnl_evolution}. This approach is
therefore recorded here as a failed attempt: it does not derive the
complementarity it sets out to model.
\end{remark}

\section{An elementary obstruction to exact canonical conjugacy}
\label{sec:obstruction}

The classical Heisenberg relation asks for operators $\hat A,\hat B$
with $[\hat A,\hat B]=ic\hat I$, $c\neq0$ a real scalar. It is standard
that such a relation cannot be realized by two \emph{bounded}
operators on a Hilbert space.

\begin{remark}[Standard background]
\label{rem:wintner_wielandt}
If $\hat A,\hat B$ are bounded operators on a Hilbert space and
$[\hat A,\hat B]=c\hat I$ for a nonzero scalar $c$, no such pair
exists; equivalently, at least one of the two operators satisfying an
exact canonical commutation relation must be unbounded
\cite{Wintner1947,Wielandt1949}.
\end{remark}

For the ordinal--M\"obius pair a stronger, elementary, and
self-contained statement holds, not requiring both operators to be
bounded.

\begin{definition}
\label{def:position_operator}
Let $S\subseteq\Z$ be a set with no repeated elements (e.g.\ $S=\N$),
and let $\{|n\rangle\}_{n\in S}$ be an orthonormal family. The
\emph{position operator} $\hat n$ is defined on the algebraic span of
this family by $\hat n|n\rangle=n|n\rangle$.
\end{definition}

\begin{proposition}
\label{prop:no_exact_conjugate}
Let $\hat n$ be as in Definition~\ref{def:position_operator}. There is
no linear operator $\hat B$, defined at least on the span of
$\{|n\rangle\}_{n\in S}$ (bounded or not), satisfying
$[\hat n,\hat B]=ic\hat I$ for a nonzero scalar $c$.
\end{proposition}

\begin{proof}
Write $B_{mn}=\langle m|\hat B|n\rangle$. Taking matrix elements of
$[\hat n,\hat B]=ic\hat I$ between $\langle m|$ and $|n\rangle$ gives
\begin{equation}
(m-n)\,B_{mn} = ic\,\delta_{mn}.
\label{eq:obstruction_matrix_elements}
\end{equation}
For $m\neq n$, since $S$ has no repeated elements $m-n\neq0$, so
\eqref{eq:obstruction_matrix_elements} forces $B_{mn}=0$: any $\hat B$
satisfying the relation is diagonal in this basis. For $m=n$, the left
side of \eqref{eq:obstruction_matrix_elements} is $0$, while the right
side is $ic\neq0$, a contradiction.
\end{proof}

\begin{corollary}
\label{cor:mu_not_conjugate}
Let $\hat\mu$ be any operator with $\hat\mu|n\rangle=\mu(n)|n\rangle$
absorbed into a basis change, or more generally any operator with
spectrum contained in $\{-1,0,1\}$. Then $\hat\mu$ cannot satisfy
$[\hat n,\hat\mu]=ic\hat I$ for any nonzero $c$. In particular, the
literal reading of Question~\ref{q:motivating} -- $\mu(n)$ as the exact
canonical momentum conjugate to $n$, in the sense $[\hat n,\hat\mu]=i\hbar\hat I$
-- is not realizable, regardless of boundedness.
\end{corollary}

\begin{proof}
Immediate from Proposition~\ref{prop:no_exact_conjugate}, which places
no boundedness hypothesis on $\hat B$.
\end{proof}

\begin{remark}
\label{rem:obstruction_consequence}
Corollary~\ref{cor:mu_not_conjugate} rules out the most literal form of
Question~\ref{q:motivating}. The two remaining sections propose two
different, weaker, but fully rigorous readings of the question, one
Hilbert-space based (Section~\ref{sec:number_phase}) and one not
(Section~\ref{sec:non_hilbert}).
\end{remark}

\section{A number-phase realization on $\ell^2(\Z)$}
\label{sec:number_phase}

\subsection{Construction}
\label{subsec:number_phase_construction}

This construction adapts the number-phase formalism of quantum optics
\cite{SusskindGlogower1964,CarruthersNieto1968} -- there developed for
the photon-number operator and its conjugate phase -- to the ordinal
variable $n$.

\begin{definition}
\label{def:shift_operator}
On $\Hil=\ell^2(\Z)$ with orthonormal basis $\{|n\rangle\}_{n\in\Z}$,
define the position operator $\hat n$ as in
Definition~\ref{def:position_operator} and the unitary shift operator
$\hat S$ by $\hat S|n\rangle=|n+1\rangle$.
\end{definition}

\begin{proposition}[Weyl relation]
\label{prop:weyl_relation}
\begin{equation}
[\hat n,\hat S]=\hat S.
\label{eq:weyl_relation}
\end{equation}
\end{proposition}

\begin{proof}
$\hat n\hat S|n\rangle=\hat n|n+1\rangle=(n+1)|n+1\rangle$, and
$\hat S\hat n|n\rangle=n|n+1\rangle$. Subtracting gives
$[\hat n,\hat S]|n\rangle=|n+1\rangle=\hat S|n\rangle$ for every basis
vector, hence \eqref{eq:weyl_relation}.
\end{proof}

\begin{remark}
\label{rem:weyl_not_obstructed}
Relation \eqref{eq:weyl_relation} does not conflict with
Proposition~\ref{prop:no_exact_conjugate}: the right-hand side is
$\hat S$, not a nonzero scalar multiple of $\hat I$.
\end{remark}

By Pontryagin duality between $\Z$ and the circle $\mathbb{T}=[-\pi,\pi)$,
write $|\theta\rangle=\sum_{n\in\Z}e^{in\theta}|n\rangle$ (formally);
then $\hat S=e^{i\hat\theta}$ in this representation, and $\hat\theta$
is the genuine conjugate variable to $\hat n$.

\begin{definition}
\label{def:coarse_mobius_operator}
Fix a measurable partition $\mathbb{T}=\Delta_{-1}\sqcup\Delta_0\sqcup\Delta_{+1}$
into three arcs. Define the bounded self-adjoint operator
\begin{equation}
\hat\mu \;=\; (-1)\,\Pi_{\Delta_{-1}} + 0\cdot\Pi_{\Delta_0} + (+1)\,\Pi_{\Delta_{+1}},
\label{eq:coarse_mobius_operator}
\end{equation}
where $\Pi_{\Delta_c}$ is the spectral projection of $\hat\theta$ onto
the arc $\Delta_c$.
\end{definition}

The partition in Definition~\ref{def:coarse_mobius_operator} is a free
choice of the construction; nothing below determines it from the
arithmetic of $\mu$.

\begin{proposition}
\label{prop:flat_distribution}
For every $n\in\Z$ and every $c\in\{-1,0,1\}$,
\begin{equation}
P(\hat\mu=c \mid \hat n = n) = \frac{|\Delta_c|}{2\pi},
\label{eq:flat_distribution}
\end{equation}
independently of $n$.
\end{proposition}

\begin{proof}
In the $\theta$-representation, $|n\rangle$ corresponds to the function
$\theta\mapsto e^{-in\theta}$, of constant modulus $1$ on $\mathbb{T}$.
Hence $P(\hat\mu=c\mid \hat n=n)=\int_{\Delta_c}|e^{-in\theta}|^2\,\frac{d\theta}{2\pi}
= |\Delta_c|/2\pi$, which does not depend on $n$.
\end{proof}

\begin{corollary}
\label{cor:no_exceptional_state}
If every $\Delta_c$ has $0<|\Delta_c|<2\pi$, then
$\operatorname{Var}(\hat\mu\mid \hat n=n)>0$ for every $n$: there is no
exceptional value of $n$ for which the outcome of $\hat\mu$ becomes
determined.
\end{corollary}

\begin{proof}
Immediate from \eqref{eq:flat_distribution}: the three probabilities
are positive constants independent of $n$, so the variance is the same
positive constant for every $n$.
\end{proof}

\begin{remark}
\label{rem:no_vacuum_analogue}
In ordinary quantum mechanics, non-commuting observables can share an
exceptional common eigenstate when the commutator's expectation value
vanishes there (e.g.\ the zero-angular-momentum state for
$\hat L_x,\hat L_y,\hat L_z$). Corollary~\ref{cor:no_exceptional_state}
shows that this pair, $(\hat n,\hat\mu)$, has no analogous exceptional
state: the incompatibility holds uniformly.
\end{remark}

\section{A non-Hilbert-space realization}
\label{sec:non_hilbert}

\subsection{Ambient space}
\label{subsec:non_hilbert_ambient}

\begin{definition}
\label{def:arithmetic_sequence_space}
Let $A=\C^\N$, the space of all sequences $a:\N\to\C$, equipped with
the topology of pointwise convergence. $A$ is a Fr\'echet space; it
carries no inner product and is not normable in general.
\end{definition}

\begin{definition}
\label{def:dirichlet_convolution_operator}
For an arithmetic function $f:\N\to\C$, the Dirichlet convolution
operator $T_f:A\to A$ is
\begin{equation}
(T_f a)(n) = \sum_{d\mid n} f(d)\,a(n/d).
\label{eq:dirichlet_convolution_operator}
\end{equation}
The set of arithmetic functions with pointwise addition and Dirichlet
convolution forms a commutative ring, isomorphic to a formal power
series ring over $\C$ in countably many indeterminates and, in
particular, a unique factorization domain \cite{CashwellEverett1959}.
\end{definition}

\begin{definition}
\label{def:pointwise_position_operator}
The position operator $\hat N:A\to A$ is $(\hat N a)(n)=n\cdot a(n)$.
\end{definition}

Note that $T_\zeta$ (convolution with the constant function
$\zeta\equiv1$) and $T_\mu$ are mutually inverse,
\begin{equation}
T_\mu T_\zeta = T_\zeta T_\mu = \hat I,
\label{eq:mobius_inversion_operator}
\end{equation}
this being the classical M\"obius inversion formula restated as an
operator identity.

\subsection{The commutator}
\label{subsec:non_hilbert_commutator}

\begin{proposition}
\label{prop:non_hilbert_commutator}
For any arithmetic function $f$,
\begin{equation}
[\hat N, T_f] = \hat N\,T_g,
\qquad
g(d) = f(d)\left(1-\frac1d\right).
\label{eq:non_hilbert_commutator}
\end{equation}
\end{proposition}

\begin{proof}
Directly from Definitions~\ref{def:dirichlet_convolution_operator} and
\ref{def:pointwise_position_operator}, for $a\in A$,
\begin{align}
(\hat N T_f a)(n) - (T_f \hat N a)(n)
&= n\sum_{d\mid n}f(d)a(n/d) - \sum_{d\mid n}f(d)\,\frac{n}{d}\,a(n/d)
\label{eq:commutator_step_one}\\
&= n\sum_{d\mid n} f(d)\left(1-\frac1d\right)a(n/d)
\label{eq:commutator_step_two}\\
&= n\,(T_g a)(n)
= (\hat N T_g a)(n),
\label{eq:commutator_step_three}
\end{align}
with $g$ as in \eqref{eq:non_hilbert_commutator}.
\end{proof}

\begin{corollary}
\label{cor:mobius_commutator_nonzero}
Taking $f=\mu$, $g(d)=\mu(d)(d-1)/d$ is not identically zero (e.g.\
$g(2)=-1/2$), so $[\hat N,T_\mu]\neq0$.
\end{corollary}

\begin{remark}[Consistency check]
\label{rem:leibniz_consistency}
Since $[\hat N,\hat I]=0$, the identity
\eqref{eq:mobius_inversion_operator} and the Leibniz rule for
commutators, $[\hat N, T_\mu T_\zeta] = [\hat N,T_\mu]T_\zeta + T_\mu[\hat N,T_\zeta]$,
give a relation between the commutators of $\hat N$ with $T_\mu$ and
with $T_\zeta$. This relation has not been developed further here.
\end{remark}

\subsection{The escape from Hilbert space}
\label{subsec:non_hilbert_escape}

\begin{definition}
\label{def:indicator_sequence}
Let $e_n\in A$ be the indicator sequence, $e_n(k)=1$ if $k=n$ and $0$
otherwise: the position eigenstate at $n$.
\end{definition}

\begin{proposition}
\label{prop:infinite_support}
\begin{equation}
(T_\mu e_n)(m) = \mu(m/n)\cdot[\,n\mid m\,].
\label{eq:infinite_support}
\end{equation}
In particular, $T_\mu e_n$ is supported on the infinite set of
multiples of $n$.
\end{proposition}

\begin{proof}
By \eqref{eq:dirichlet_convolution_operator}, $(T_\mu e_n)(m)=\sum_{d\mid m}\mu(d)e_n(m/d)$.
The summand is nonzero only when $m/d=n$, i.e.\ $d=m/n$, which requires
$n\mid m$; when this holds, the unique contributing term is $\mu(m/n)$.
\end{proof}

\begin{corollary}
\label{cor:no_hilbert_space}
$T_\mu e_n\notin\ell^p(\N)$ for any $1\le p<\infty$: applying the
momentum-type operator $T_\mu$ to an exact position eigenstate
produces a vector that cannot be normalized in any Hilbert (or, more
generally, $\ell^p$) space, and exists only in the Fr\'echet space $A$
of Definition~\ref{def:arithmetic_sequence_space}.
\end{corollary}

\begin{proof}
By Proposition~\ref{prop:infinite_support}, $|(T_\mu e_n)(m)|\in\{0,1\}$
on the multiples of $n$, an infinite set on which the entries do not
tend to $0$; hence the sequence is not in $\ell^p$ for any finite $p$.
\end{proof}

\begin{remark}
\label{rem:non_hilbert_motivated}
Corollary~\ref{cor:no_hilbert_space} is the sense in which the
non-Hilbert ambient space of
Definition~\ref{def:arithmetic_sequence_space} is not a stylistic
choice: it is forced by the requirement that $\hat N$ (position, exact
ordinal) and $T_\mu$ (the Dirichlet-convolution momentum-type
operator) act on a common domain. No probability distribution is
invoked anywhere in this section; the incompatibility recorded in
Corollary~\ref{cor:no_hilbert_space} is a deterministic algebraic fact.
\end{remark}

\section{A Bost--Connes-type flow on the Fr\'echet space}
\label{sec:bc_flow}

This section pushes the construction of Section~\ref{sec:non_hilbert}
in the direction of the Bost--Connes quantum statistical mechanical
system \cite{BostConnes1995}, which is built on a genuine Hilbert
space. The goal here is more modest: to see how much of that structure
-- a Hamiltonian, a time evolution, and a distinguished commutator --
survives when the ambient space is the Fr\'echet space $A$ of
Definition~\ref{def:arithmetic_sequence_space} instead.

\subsection{The ring homomorphism property}
\label{subsec:bc_flow_ring_hom}

\begin{proposition}
\label{prop:convolution_operator_homomorphism}
For arithmetic functions $f,g$,
\begin{equation}
T_f T_g = T_{f*g},
\label{eq:convolution_operator_homomorphism}
\end{equation}
where $*$ denotes Dirichlet convolution. In particular, since $*$ is
commutative, the operators $\{T_f\}_f$ pairwise commute.
\end{proposition}

\begin{proof}
For $a\in A$,
\begin{align}
(T_fT_g a)(n)
&= \sum_{d\mid n} f(d)\,(T_g a)(n/d)
= \sum_{d\mid n} f(d) \sum_{e\mid (n/d)} g(e)\,a\!\left(\frac{n}{de}\right)
\label{eq:hom_step_one}\\
&= \sum_{de\mid n} f(d)g(e)\,a\!\left(\frac{n}{de}\right)
= \sum_{k\mid n}\Big(\sum_{d\mid k} f(d)g(k/d)\Big) a(n/k)
\label{eq:hom_step_two}\\
&= \sum_{k\mid n} (f*g)(k)\,a(n/k)
= (T_{f*g}a)(n).
\label{eq:hom_step_three}
\end{align}
\end{proof}

\begin{remark}
\label{rem:commutative_subalgebra}
Proposition~\ref{prop:convolution_operator_homomorphism} shows that
$\{T_f\}_f$ is a commutative subalgebra of operators on $A$: the
non-commutativity found in
Proposition~\ref{prop:non_hilbert_commutator} comes entirely from the
position/Hamiltonian-type operator, not from within this family. This
mirrors the structure of the Bost--Connes algebra, where the Hecke
operators also form a commutative subalgebra and non-commutativity
enters only through the number operator \cite{BostConnes1995}.
\end{remark}

\subsection{Hamiltonian and time evolution}
\label{subsec:bc_flow_hamiltonian}

\begin{definition}
\label{def:hamiltonian_operator}
Define $\hat H:A\to A$ by $(\hat H a)(n) = \log(n)\,a(n)$ for $n\ge2$,
and $(\hat Ha)(1)=0$.
\end{definition}

\begin{proposition}
\label{prop:time_evolution}
For $t\in\mathbb{R}$, define formally $\sigma_t(T_f) := e^{it\hat H}T_f e^{-it\hat H}$.
Then
\begin{equation}
\sigma_t(T_f) = T_{f_t}, \qquad f_t(d) = f(d)\,d^{it}.
\label{eq:time_evolution}
\end{equation}
\end{proposition}

\begin{proof}
$(e^{it\hat H}T_f e^{-it\hat H}a)(n) = n^{it}\sum_{d\mid n} f(d)\,(e^{-it\hat H}a)(n/d)
= n^{it}\sum_{d\mid n} f(d)(n/d)^{-it}a(n/d) = \sum_{d\mid n}f(d)d^{it}a(n/d) = (T_{f_t}a)(n)$.
\end{proof}

\begin{remark}[Parallel with Bost--Connes]
\label{rem:bc_parallel}
The twist $f(d)\mapsto f(d)d^{it}$ in \eqref{eq:time_evolution} is the
same qualitative operation as the time evolution of the Bost--Connes
system, which acts on its generating isometries $\mu_n$ by
$\sigma_t(\mu_n)=n^{it}\mu_n$ \cite{BostConnes1995}. This is recorded
as an adaptation [ADAPT]: same mechanism (multiplicative twist by
$d^{it}$, generated by a logarithmic Hamiltonian), transplanted from a
$C^*$-algebra of isometries on $\ell^2(\N)$ to the commutative
convolution algebra $\{T_f\}$ on the Fr\'echet space $A$. No claim is
made that $\sigma_t$ as defined here is continuous in any topology on
$A$, or that it extends to a genuine one-parameter group of
automorphisms of a well-defined algebra; this is left as
Question~\ref{q:flow_continuity} below.
\end{remark}

\subsection{The commutator with the Hamiltonian}
\label{subsec:bc_flow_commutator}

\begin{proposition}
\label{prop:hamiltonian_commutator}
For any arithmetic function $f$,
\begin{equation}
[\hat H, T_f] = T_{f\log}, \qquad (f\log)(d) := f(d)\log(d).
\label{eq:hamiltonian_commutator}
\end{equation}
\end{proposition}

\begin{proof}
For $a\in A$ and $n\ge1$,
\begin{align}
(\hat H T_f a)(n) - (T_f\hat H a)(n)
&= \log(n)\sum_{d\mid n} f(d)a(n/d) - \sum_{d\mid n} f(d)\log(n/d)\,a(n/d)
\label{eq:ham_comm_step_one}\\
&= \sum_{d\mid n} f(d)\big[\log n - \log(n/d)\big]a(n/d)
= \sum_{d\mid n} f(d)\log(d)\,a(n/d)
\label{eq:ham_comm_step_two}\\
&= (T_{f\log}a)(n).
\label{eq:ham_comm_step_three}
\end{align}
\end{proof}

\begin{corollary}
\label{cor:hamiltonian_mobius_commutator}
\begin{equation}
[\hat H, T_\mu] = -\,T_\Lambda\, T_\mu,
\label{eq:hamiltonian_mobius_commutator}
\end{equation}
where $\Lambda$ is the von Mangoldt function.
\end{corollary}

\begin{proof}
The classical identity $\Lambda = \mu * \log$, i.e.
$\Lambda(n)=\sum_{d\mid n}\mu(d)\log(n/d)$, is standard
\cite[Thm.~2.10]{Apostol1976}. Expanding $\log(n/d)=\log n-\log d$ and
using $\sum_{d\mid n}\mu(d)=\varepsilon(n)$ (the Dirichlet identity
element) gives, for $n\ge1$,
\begin{equation}
(\mu\log)*\zeta = -\Lambda,
\label{eq:mulog_convolution_identity}
\end{equation}
where $(\mu\log)(d)=\mu(d)\log(d)$. Since $\zeta^{-1}=\mu$ in the
Dirichlet ring, \eqref{eq:mulog_convolution_identity} gives
$\mu\log = -\Lambda*\mu$. By
Proposition~\ref{prop:convolution_operator_homomorphism},
$T_{\mu\log}=T_{-\Lambda*\mu}=-T_\Lambda T_\mu$. Combined with
Proposition~\ref{prop:hamiltonian_commutator} at $f=\mu$, this gives
\eqref{eq:hamiltonian_mobius_commutator}.
\end{proof}

\begin{remark}
\label{rem:hamiltonian_mobius_meaning}
Corollary~\ref{cor:hamiltonian_mobius_commutator} refines the bare
non-vanishing statement of
Corollary~\ref{cor:mobius_commutator_nonzero}: the commutator of the
Hamiltonian $\hat H$ with the M\"obius operator $T_\mu$ is not an
arbitrary nonzero operator but is exactly (minus) the von Mangoldt
operator composed with $T_\mu$ itself. The incompatibility between
$\hat H$ (energy/ordinal scale) and $T_\mu$ (M\"obius momentum) is thus
governed, in this formalism, by the same function $\Lambda$ that
records the location and multiplicity of prime powers -- a fact
obtained here as a derived consequence of
Definitions~\ref{def:dirichlet_convolution_operator} and
\ref{def:hamiltonian_operator}, not assumed.
\end{remark}

\subsection{A formal partition function}
\label{subsec:bc_flow_partition}

\begin{heuristic}
\label{heur:formal_partition_function}
On $\ell^2(\N)$, the Bost--Connes Hamiltonian has partition function
$\operatorname{Tr}(e^{-\beta H}) = \sum_n n^{-\beta} = \zeta(\beta)$, a
rigorous trace of a well-defined trace-class operator for
$\operatorname{Re}\beta>1$ \cite{BostConnes1995}. On the Fr\'echet
space $A$ there is no inner product and no trace-class structure, so
no analogous trace exists rigorously here. Formally associating to the
diagonal operator $e^{-\beta\hat H}$ (with eigenvalues $n^{-\beta}$ on
$e_n$) the generating Dirichlet series $\sum_n n^{-\beta}=\zeta(\beta)$
reproduces the same expression, but this is a bookkeeping device, not
a trace of an operator on $A$ in any sense established in this
notebook. This is recorded as a heuristic only.
\end{heuristic}

\begin{question}
\label{q:flow_continuity}
Is there a topology on (a subspace of) $A$, coarser than or different
from the Fr\'echet topology of
Definition~\ref{def:arithmetic_sequence_space}, in which $\sigma_t$ of
Proposition~\ref{prop:time_evolution} is continuous in $t$ and the
operators $T_f$ are continuous, so that
Heuristic~\ref{heur:formal_partition_function} could be replaced by an
actual trace or an actual KMS condition?
\end{question}

\section{Related structures in the literature}
\label{sec:literature}

This section records what has and has not been located in the
literature, and is not itself a source of new results.

\begin{itemize}[leftmargin=*]

\item The classical probabilistic model for $\mu$ replaces it by a
Rademacher random multiplicative function $f$, $f(p)=\pm1$ i.i.d.\ on
primes \cite{Wintner1944}. A na\"ive central-limit heuristic for
$\sum_{n\le x}f(n)$, of the same flavour as the ``fluctuating version''
discarded in Remark~\ref{rem:failed_mqmnl}, predicts fluctuations of
order $\sqrt{x\log\log x}$. This has been shown to be false in the
opposite direction: $\mathbb{E}\left|\sum_{n\le x}f(n)\right|\asymp
\sqrt{x}/(\log\log x)^{1/4}$ \cite{Harper2020}, i.e.\ \emph{better}
than square-root cancellation, driven by multiplicative correlations
among the $f(n)$ rather than by any added source of uncertainty. This
is the opposite of what an add-on uncertainty postulate such as
\eqref{eq:mqmnl_postulated_uncertainty} would suggest.

\item The Hilbert--P\'olya program and the Berry--Keating $H=xp$
model \cite{BerryKeating1999} pursue a different question -- realizing
the ordinates of the nontrivial zeros of $\zeta(s)$ as the spectrum of
a self-adjoint operator -- and are not directly about $\mu(n)$ or about
the ordinal--M\"obius pairing studied here.

\item The construction of Section~\ref{sec:non_hilbert} is structurally
closest to the Bost--Connes system \cite{BostConnes1995}, whose scaling
isometries $\mu_n$ and phase-type generators $e(r)$ satisfy a
non-commutation relation of the same qualitative type as
\eqref{eq:non_hilbert_commutator} (multiplication/averaging over
multiples of $n$ against a phase-type variable). That system is built
on a genuine Hilbert space and a $C^*$-algebra of isometries, motivated
by class field theory and KMS states, and does not use the
Fr\'echet-space escape argument of
Corollary~\ref{cor:no_hilbert_space}.

\end{itemize}

\begin{remark}[Novelty check]
\label{rem:novelty_check}
A literature search was carried out for the specific commutator
identity \eqref{eq:non_hilbert_commutator} and for the escape argument
of Corollary~\ref{cor:no_hilbert_space}. Neither was located in this
exact form. Absence of a search result is weak evidence, not a proof
of originality; the pieces used (the Dirichlet convolution ring
\cite{CashwellEverett1959}, elementary computation) are all standard.
\end{remark}

\section{Open questions}
\label{sec:open}

\begin{question}
\label{q:relate_to_bjorner}
Does the commutator identity \eqref{eq:non_hilbert_commutator}, or the
escape argument of Corollary~\ref{cor:no_hilbert_space}, connect in any
precise way to the Bj\"orner cellular complex realization of $M(n)$ as
a reduced Euler characteristic, used elsewhere in this research
program?
\end{question}

\begin{question}
\label{q:other_partitions}
How does Corollary~\ref{cor:no_exceptional_state} depend on the choice
of partition $\Delta_{-1},\Delta_0,\Delta_{+1}$ in
Definition~\ref{def:coarse_mobius_operator}? Is there a
partition-independent formulation of the same complementarity
statement?
\end{question}

\begin{conjecture}[Purely exploratory]
\label{conj:leibniz_extension}
The relation recorded in Remark~\ref{rem:leibniz_consistency} between
$[\hat N,T_\mu]$ and $[\hat N,T_\zeta]$ extends to a full family of
commutator identities $[\hat N,T_f]$ for $f$ ranging over the
multiplicative functions, organized by the Dirichlet-series structure
of $f$. No evidence beyond Proposition~\ref{prop:non_hilbert_commutator}
is offered for this at present.
\end{conjecture}

\section*{Notebook History}
\addcontentsline{toc}{section}{Notebook History}

\noindent Version 0.1: Naive quantum ansatz (MQMNL) recorded and
identified as a failed attempt (Section~\ref{sec:naive_ansatz}).

\noindent Version 0.2: Elementary obstruction to exact canonical
conjugacy established (Section~\ref{sec:obstruction}).

\noindent Version 0.3: Number-phase Hilbert-space construction added
(Section~\ref{sec:number_phase}).

\noindent Version 0.4: Non-Hilbert Dirichlet-convolution construction
added, including the escape-from-Hilbert-space argument
(Section~\ref{sec:non_hilbert}).

\noindent Version 0.5: Bost--Connes-type Hamiltonian, time evolution,
and commutator with $T_\mu$ added, expressing the incompatibility
between $\hat H$ and $T_\mu$ exactly in terms of the von Mangoldt
operator; formal partition function flagged as heuristic
(Section~\ref{sec:bc_flow}).

\begin{thebibliography}{99}

\bibitem{Wintner1947}
A.~Wintner, \emph{The unboundedness of quantum-mechanical matrices},
Phys. Rev. \textbf{71} (1947), 738--739.
doi:10.1103/PhysRev.71.738.2

\bibitem{Wielandt1949}
H.~Wielandt, \emph{\"Uber die Unbeschr\"anktheit der Operatoren der
Quantenmechanik}, Math. Ann. \textbf{121} (1949), 21.
doi:10.1007/BF01329611

\bibitem{SusskindGlogower1964}
L.~Susskind and J.~Glogower, \emph{Quantum mechanical phase and time
operator}, Physics \textbf{1} (1964), 49--61.

\bibitem{CarruthersNieto1968}
P.~Carruthers and M.~M.~Nieto, \emph{Phase and angle variables in
quantum mechanics}, Rev. Mod. Phys. \textbf{40} (1968), 411--440.
doi:10.1103/RevModPhys.40.411

\bibitem{CashwellEverett1959}
E.~D.~Cashwell and C.~J.~Everett, \emph{The ring of number-theoretic
functions}, Pacific J. Math. \textbf{9} (1959), no.~4, 975--985.
doi:10.2140/pjm.1959.9.975

\bibitem{Wintner1944}
A.~Wintner, \emph{Random factorisations and Riemann's hypothesis},
Duke Math. J. \textbf{11} (1944), 267--275.

\bibitem{Harper2020}
A.~J.~Harper, \emph{Moments of random multiplicative functions, I: Low
moments, better than squareroot cancellation, and critical
multiplicative chaos}, Forum Math. Pi \textbf{8} (2020), e1, 95~pp.
arXiv:1703.06654

\bibitem{BostConnes1995}
J.-B.~Bost and A.~Connes, \emph{Hecke algebras, type III factors and
phase transitions with spontaneous symmetry breaking in number
theory}, Selecta Math. (N.S.) \textbf{1} (1995), 411--457.

\bibitem{Apostol1976}
T.~M.~Apostol, \emph{Introduction to Analytic Number Theory},
Undergraduate Texts in Mathematics, Springer-Verlag, New York, 1976.

\bibitem{BerryKeating1999}
M.~V.~Berry and J.~P.~Keating, \emph{The Riemann zeros and eigenvalue
asymptotics}, SIAM Rev. \textbf{41} (1999), 236--266.

\end{thebibliography}

\end{document}
